\phantomsection
\addcontentsline{toc}{chapter}{课程大纲与课历}
\markboth{课程大纲与课历}{}
\noindent{\color{structurecolor}\bfseries\fontsize{14.4}{20}\selectfont 课程大纲与课历}\par

\textbf{课程与授课安排.} MIT 数学系本科课程 18.330, Introduction to Numerical Analysis, 2012 年春季, 教师 Laurent Demanet. 每周两次讲课, 每次 1.5 小时. 官方公开资料没有给出每次讲课的具体日期、钟点或教室; 下表仅列官方讲义目录的讲次范围, 不将作业截止日当作授课日.

\textbf{先修与编程.} 先修 18.01 微积分、18.02 多元微积分、18.03 微分方程. 18.06 水平的矩阵与线性代数知识有帮助, 但不强制要求. 作业涉及基本编程, 可自行选择语言; 原课程推荐 MATLAB, 并鼓励尚不会的学生学习. 本整理本的复现实验使用 Python, 同时在题目明确要求时保留相应 MATLAB 核心代码.

\textbf{课程目标与大纲.} 研究如何在计算机中以数列处理函数、导数、积分及微分方程, 重点理解 Taylor、Fourier 等展开的收敛速度. 官方大纲包括: 级数展开; 积分求和与导数差分; 插值、样条及数值微积分; 常微分方程初值问题; 求根、Newton 方法与边值问题; Fourier 变换、Fourier 级数及 Shannon 采样理论; 带限插值与谱方法; 最小二乘逼近; 主成分分析. 原课程不覆盖 18.303 线性偏微分方程的内容, 线性代数的比重也明显低于 18.06SC. 七章现有讲义没有独立的最小二乘或主成分分析章节, 本次不凭大纲虚构其原始讲义.

\begin{center}
\begin{tabular}{cl}
\toprule
官方讲次 & 讲义主题 \\
\midrule
1--2 & 级数与数列 \\
3--4 & 积分求和与导数差分 \\
5--8 & 插值 \\
9--10 & 非线性方程 \\
11--15 & 常微分方程方法 \\
16--20 & Fourier 分析 \\
21--25 & 谱插值、谱微分与谱求积 \\
\bottomrule
\end{tabular}
\end{center}

\textbf{考核规则.} 作业 50\%, 课堂期中考试 20\%, 课堂期末考试 30\%. 作业包括理论题与数值实验, 不接受迟交, 去掉最低一次成绩. 允许合作讨论, 但提交的程序与书面作业必须由学生本人原创完成. 期中与期末均开卷, 不准使用计算器、手机或计算机. 官方公开页面没有考试卷、考试时长或日期; 本册不补造考试题或官方答案.

\textbf{作业日期与范围.} 八份作业共 33 道原编号题. 含附加问共冻结 68 个末级题目单元; 无子问的题记为一单元, 同条中的多个要求均保留. PS8 第 4 题末条同时包含微分与求积两项. 具体清单见源码中的 \nolinkurl{assessment-inventory.json}.
\begin{center}
\begin{tabular}{clrr}
\toprule
作业 & 原件截止信息 & 原编号题数 & 单元数 \\
\midrule
1 & 2 月 23 日 & 4 & 8 \\
2 & 3 月 1 日 & 5 & 9 \\
3 & 3 月 15 日 & 4 & 7 \\
4 & 4 月 3 日, Tuesday & 5 & 9 \\
5 & 4 月 12 日, 原印 Tuesday & 3 & 11 \\
6 & 4 月 26 日, Thursday & 3 & 9 \\
7 & 5 月 3 日, Thursday & 4 & 8 \\
8 & Not due; Fourier exercises (05/15) & 5 & 7 \\
\bottomrule
\end{tabular}
\end{center}
PS5 的原件星期与 2012 年 4 月 12 日的实际星期四不符, 上表保留原印内容, 不另推测截止日. PS8 是不需提交的 Fourier 练习, 原件没有分值. 其余各题分值逐题保留, 不自行分配子问分值.

\textbf{题面与解答身份.} 第 8--15 章分别收入 PS1--PS8 的完整中文题面、原题号及分值. \textbf{全部解答均由 AI 编写, 不是 MIT 或教师提供的官方解答.} 原件的提示和附注归入题面; 针对排印问题、缺条件或离散约定的说明标为“整理注”. 数值输出来自随源码提供的实际脚本, 不以实验代替证明.

\textbf{来源.} \href{https://ocw.mit.edu/courses/18-330-introduction-to-numerical-analysis-spring-2012/pages/syllabus/}{官方 Syllabus}, \href{https://ocw.mit.edu/courses/18-330-introduction-to-numerical-analysis-spring-2012/pages/lecture-notes/}{官方讲义目录}, \href{https://ocw.mit.edu/courses/18-330-introduction-to-numerical-analysis-spring-2012/pages/assignments/}{官方 Assignments}. 原始八份作业及网页快照独立归档于 \nolinkurl{sources/18.330-spring-2012/}; 下载直链与 SHA-256 见 \nolinkurl{assessment-sources.json}. 题面与本改编沿用 MIT OCW 的 CC BY-NC-SA 4.0 默认许可, 个别第三方权利声明优先, 不表示 MIT 对本解答的认可.
